\section{Models}
\label{sec:models}

\subsection{Static arbitrage in total variance}

With $c(k,T)=C/(DF)$ the forward-normalized call price, a surface is free of static arbitrage if
$c$ is nondecreasing in $T$ at fixed $k$ (no calendar arbitrage) and convex and decreasing in the
strike (no butterfly or vertical arbitrage), with the usual price bounds~\citep{carr2005}. In total
variance these become~\citep{gatheral2014}
\begin{equation}
  \partial_T w(k,T) \ge 0, \qquad
  g(k) = \Bigl(1-\frac{k w'}{2w}\Bigr)^{2} - \frac{w'^{2}}{4}\Bigl(\frac{1}{w}+\frac14\Bigr)
  + \frac{w''}{2} \ge 0,
  \label{eq:conditions}
\end{equation}
where primes are derivatives in $k$ at fixed $T$: the risk-neutral density at strike $K$ is
$g(k)\,\phi(d_2)/(K\sqrt{w})$, so it is nonnegative exactly where $g$ is (\cref{app:arbitrage}). The
wings must also satisfy Lee's moment bound $\limsup_{|k|\to\infty} w/|k|\le 2$~\citep{lee2004}.

\subsection{Baselines and the penalized network}

The baselines are SVI fitted to each expiry, a power-law SSVI surface~\citep{gatheral2014}, a
smoothing spline in $k$ per expiry with total variance linear in $T$ between expiries, and an
unpenalized MLP. Only SSVI is arbitrage-free by construction.

The penalized network (\pnn) multiplies a same-day SSVI fit by a bounded correction,
$w=w_{\text{SSVI}}\exp\bigl(3\tanh(f/3)\bigr)$, where $f$ is an MLP (three hidden layers of 64
$\tanh$ units) of $k/\sqrt{\theta_{\text{SSVI}}(T)}$ and $\log T$. Its loss is the mean squared
residual in bid-ask half-spreads plus hinge penalties, each with weight 100, on the calendar and
butterfly conditions of \cref{eq:conditions} and on the wing slope $|\partial_k w|\le 2$. The
penalties are evaluated by automatic differentiation on a collocation grid that covers every
quoted maturity, the midpoints between them and strikes well beyond the quoted range, because
penalizing only at quotes leaves the surface free to violate the conditions between and beyond
them. Fits run in double precision, with 500 Adam~\citep{kingma2015} steps followed by
L-BFGS~\citep{liu1989}, warm-started from the previous trading day.

\subsection{Hard-constrained eSSVI mixtures}

An eSSVI slice~\citep{hendriks2019} with ATM total variance $\theta$, skew scale $\psi$ and
correlation $\rho$ is
\begin{equation}
  w(k)=\tfrac12\Bigl(\theta+\rho\psi k+\sqrt{(\psi k+\theta\rho)^2+\theta^2(1-\rho^2)}\Bigr).
\end{equation}
In the wing slopes $u_\pm=\psi(1\pm\rho)$, which give $w\sim u_\pm|k|/2$ as $k\to\pm\infty$, the
known sufficient conditions read: for $T_1<T_2$, $\theta_2>\theta_1$, $u_\pm(T_2)\ge u_\pm(T_1)$ and
$\psi_2/\theta_2\le\psi_1/\theta_1$ (no calendar arbitrage), and for each slice
$\max(u_+,u_-)<4$ and $\psi\max(u_+,u_-)\le 4\theta$ (no butterfly arbitrage)~\citep{gatheral2014,hendriks2019,mingone2022}.
Every condition is a monotonicity in $T$, a bound on the wings, or a lower bound on $\theta$ given
the wings. Knots are therefore generated in maturity order from unconstrained rates: the wings
first, nondecreasing and below 4, then $\theta$ above its lower bounds (\cref{app:essvi}). Between
knots $(\theta,u_+,u_-)$ are linear in $T$, which preserves the conditions~\citep{corbetta2019}.
Every parameter vector is feasible, so violations are zero by theorem rather than by tuning.

A single eSSVI surface is too rigid for SPX smiles (\cref{tab:models}). The headline model,
\mixfour, mixes four surfaces in forward-normalized price space,
$c(k,T)=\sum_{j=1}^{4}\pi_j\,c_{\text{B}}\bigl(k,w_j(k,T)\bigr)$, with weights $\pi_j\ge0$ that do not
depend on $T$; $c_{\text{B}}$ is Black's normalized call price. A convex combination of
arbitrage-free price surfaces with maturity-independent weights is arbitrage-free
(\cref{prop:mixture}). The model's total variance is the Black implied variance of the mixture price,
found by a bracketed Newton solve in log price, and fits differentiate through it implicitly. Knots
are free at each quoted maturity (\mixfour) or generated by a network of $\log T$ (\texttt{nmix4}).
The fit minimizes the squared IV residual in half-spreads with L-BFGS. Its default start is a single
eSSVI fit, whose ATM variance is scaled to seed the four components; \cref{sec:refits} adds a second
start from the previous day's surface.

\subsection{Cross-day conditional model}

One network serves every day. A set encoder reads any subset of a day's quotes, each described by
$k/\sqrt{T}$, $\log T$, its IV and its log half-spread, through a shared per-quote network with
mean and max pooling, and returns a latent $z\in\mathbb{R}^8$. The decoder maps $z$ to the rates of
four eSSVI surfaces on a fixed grid of 26 maturities (5 to 900 days) and their mixture weights,
assembled exactly as \mixfour. Every latent therefore decodes to an arbitrage-free surface,
including latents that were never encoded from a chain, such as forecasts. Training follows neural
processes~\citep{garnelo2018}: each step encodes a random partial chain (wings cut, a random share of
interior strikes dropped, sometimes whole maturities) and scores the decoded surface on all of that
day's quotes. A day is then fitted either by the encoder alone (amortized) or by refining its latent
with L-BFGS against the quotes (refined). The latent size was chosen on the early-stopping loss
(squared half-spreads 14.8, 7.6 and 8.1 for $d=4$, 8 and 16).

The SPX model trains on 2018 to September 2021, with October--December 2021 for early stopping, and
a second version on 2010 to September 2021. Walk-forward versions fine-tune a copy before each
quarter on every screened day up to its start, at a constant learning rate and for an epoch count
chosen on 2022 (40 for the 2018 model, 10 for the 2010 model). The synthetic model trains on Heston
2020, and a second version on nine generator years.

For forecasting, a VAR(1) fitted to the latents of the training days forecasts $z$ $h$ days ahead,
and the decoded surface is scored on every quote of the target day against two random walks: the
latent of the origin day, and the origin day's \mixfour{} fit.
